\section{Conclusion and future work}
\label{sec:conclusion}

A user who tries a configuration, judges the answer and resends the prompt to
a stronger one pays for every attempt, and none of the routing studies we
located measures that cost. We defined it as the escalation tax, split it
into escalation and stickiness, and measured it on three pools: 1{,}200
gradable prompts run through all eight cells of a Qwen3 size-by-mode grid on
a laptop GPU with metered board energy, thirteen public models priced by
their providers, and six hosted cells in which reasoning effort is a paid
setting.

The tax is large, and most of it comes from where a user stays rather than
how they climb. Leaving the 8B thinking cell selected pays 12.96 times the
minimum sufficient energy and answers 79.8\% of prompts, while a frugal climb
pays 0.89 times and answers 92.2\%; the same ordering holds on the public
pool, at 7.96 against 2.69 times, and on the hosted ladder, at 4.35 against
0.22. On the gradable families thinking was never the cheaper way to buy
accuracy: below 8B, the next size up without thinking answered more prompts
for 3.8 to 5.5 times less energy per solved prompt. Open-ended requests,
graded by a judge that agrees with a human only moderately, differed---
thinking cost about twice the energy per solved prompt and the largest
thinking cell was the most accurate---yet the frugal climb still paid the
least tax. Routing one tier too low cost 2.71 times as much energy as routing
one tier too high, an asymmetry that reverses when counted in tokens, as does
the ranking of the two costly habits: cost has to be measured in the unit
that is actually paid.

Whether a router beat the frugal climb turned on the recovery rule. When a
failed route was resent to the 8B thinking cell, the best router, logistic
regression on TF-IDF features, spent 1.30 times the climb's energy per solved
prompt and answered 3.7 points fewer; when failures climbed from the routed
cell instead, every router came within 6\% of it and four were 2.6\% to 3.4\%
cheaper. Against the habits that start high the router spent 64\% less than
leaving the 8B thinking cell selected and 37\% less than a self-verifying
cascade, both per solved prompt, and where users accept 15\% of wrong answers
it answered 7.1 points more than the climb. Charging routers for their own
energy barely changed their ranking, since none cost more than 7.6\,J per
decision. For a user without a router, the evidence points to a short rule:
start small with reasoning off, move up a size before turning reasoning on,
and treat extended reasoning as a choice made per prompt rather than a
setting left on.

Several directions follow. The constructed policies and swept detection
probabilities are the least grounded part of the comparator, and a study that
records how people actually move between the configurations of a tiered
service would replace those assumptions with estimates. Treating the
reasoning budget as a continuous axis rather than a switch would turn the
grid into a surface and the router's decision into a choice of size and
budget; the results make this the most pressing direction, since the cap is a
budget fixed once for every prompt. Because energy per output token barely
changes with the mode within a size while changing by a factor of 6.6 between
sizes, a table of joules per output token per model, measured once, could
turn the token logs services already keep into energy estimates. A router
could also be trained against the recovery rule it will face, since the
threshold that minimises cost when failures jump to $c_8$ is not the one that
minimises it when they climb, or could propose a starting rung for the frugal
climb instead of a final cell. Finally, a cost model in which a retry carries
the conversation so far, and in which changing size pays for loading the new
model, would bring the tax closer to what users of memory-limited hardware
pay.
